
\documentclass[12pt,a4paper]{article}

% --------------------------------------------------
% Packages
% --------------------------------------------------
\usepackage{geometry}
\usepackage{graphicx}
\usepackage{amsmath}
\usepackage{array}
\usepackage{booktabs}
\usepackage{listings}
\usepackage{hyperref}
\usepackage{setspace}
\usepackage{fontspec}
\usepackage{xepersian}

% --------------------------------------------------
% Page layout
% --------------------------------------------------
\geometry{
    top=2.5cm,
    bottom=2.5cm,
    left=2.5cm,
    right=2.5cm
}

\onehalfspacing
\setlength{\parskip}{6pt}
\pagestyle{plain}

% --------------------------------------------------
% Cross-platform fonts
% --------------------------------------------------
\IfFontExistsTF{Amiri}{
    \settextfont{Amiri}
}{
    \IfFontExistsTF{Tahoma}{
        \settextfont{Tahoma}
    }{
        \settextfont{DejaVu Sans}
    }
}

\IfFontExistsTF{Latin Modern Roman}{
    \setlatintextfont{Latin Modern Roman}
}{
    \setlatintextfont{Times New Roman}
}

% --------------------------------------------------
% Code formatting
% --------------------------------------------------
\lstset{
    basicstyle=\ttfamily\small,
    breaklines=true,
    frame=single,
    numbers=left,
    numberstyle=\tiny,
    showstringspaces=false
}

% --------------------------------------------------
% Report information
% Edit these fields for your own report
% --------------------------------------------------
\newcommand{\universityname}{دانشگاه نمونه}
\newcommand{\coursename}{آزمایشگاه سیستم‌عامل}
\newcommand{\reporttitle}{آشنایی با محیط لینوکس}
\newcommand{\experimentnumber}{آزمایش اول}
\newcommand{\studentname}{نام و نام خانوادگی}
\newcommand{\studentid}{۱۴۰۳۰۰۰۰۰۰۰}
\newcommand{\professorname}{نام استاد}
\newcommand{\submissiondate}{۱۴۰۵/۰۷/۲۷}
\newcommand{\academicyear}{۱۴۰۵--۱۴۰۶}

\begin{document}

% --------------------------------------------------
% Title page
% --------------------------------------------------
\begin{titlepage}
\begin{center}

\vspace*{1cm}

{\Large \universityname}

\vspace{1.2cm}

{\Huge\bfseries \coursename}

\vspace{0.8cm}

{\LARGE \experimentnumber}

\vspace{0.4cm}

{\Large \reporttitle}

\vspace{2cm}

{\large نام دانشجو: \studentname}

\vspace{0.5cm}

{\large شماره دانشجویی: \studentid}

\vspace{0.5cm}

{\large استاد درس: \professorname}

\vspace{0.5cm}

{\large تاریخ تحویل: \submissiondate}

\vfill

{\large سال تحصیلی \academicyear}

\end{center}
\end{titlepage}

% --------------------------------------------------
% Table of contents
% --------------------------------------------------
\tableofcontents
\clearpage

% --------------------------------------------------
% Introduction
% --------------------------------------------------
\section{مقدمه}

سیستم‌عامل یکی از اجزای اصلی هر سیستم رایانه‌ای است
که وظیفه مدیریت منابع سخت‌افزاری و فراهم کردن
محیط مناسب برای اجرای برنامه‌ها را بر عهده دارد.

در این آزمایش با محیط سیستم‌عامل
\lr{Ubuntu}
و تعدادی از دستورات پایه ترمینال آشنا می‌شویم.

% --------------------------------------------------
% Objectives
% --------------------------------------------------
\section{هدف آزمایش}

اهداف این آزمایش عبارت‌اند از:

\begin{enumerate}
    \item آشنایی با محیط سیستم‌عامل لینوکس
    \item یادگیری نحوه استفاده از ترمینال
    \item اجرای دستورات پایه
    \item ثبت و تحلیل نتایج آزمایش
    \item مستندسازی مراحل با استفاده از لاتک
\end{enumerate}

% --------------------------------------------------
% Tools
% --------------------------------------------------
\section{ابزارهای موردنیاز}

\begin{itemize}
    \item سیستم‌عامل \lr{Ubuntu}
    \item نرم‌افزار ترمینال
    \item ویرایشگر \lr{TeXstudio}
    \item موتور پردازش \lr{XeLaTeX}
\end{itemize}

% --------------------------------------------------
% Experiment steps
% --------------------------------------------------
\section{مراحل انجام آزمایش}

\subsection{بررسی مسیر فعلی}

ابتدا ترمینال را باز می‌کنیم و دستور زیر را اجرا می‌کنیم:

\begin{latin}
\begin{lstlisting}
pwd
\end{lstlisting}
\end{latin}

این دستور مسیر پوشه فعلی را نمایش می‌دهد.

\subsection{نمایش فایل‌ها و پوشه‌ها}

برای مشاهده محتوای پوشه فعلی از دستور زیر
استفاده می‌کنیم:

\begin{latin}
\begin{lstlisting}
ls -la
\end{lstlisting}
\end{latin}

گزینه‌های \lr{-l} و \lr{-a} به‌ترتیب اطلاعات
بیشتر و فایل‌های مخفی را نمایش می‌دهند.

\subsection{بررسی نام کاربر}

دستور زیر نام کاربر فعال را نمایش می‌دهد:

\begin{latin}
\begin{lstlisting}
whoami
\end{lstlisting}
\end{latin}

\subsection{بررسی اطلاعات هسته}

برای نمایش اطلاعات هسته و سیستم از دستور زیر
استفاده می‌کنیم:

\begin{latin}
\begin{lstlisting}
uname -a
\end{lstlisting}
\end{latin}

\subsection{بررسی نسخه اوبونتو}

برای مشاهده اطلاعات نسخه توزیع می‌توانیم
از دستور زیر استفاده کنیم:

\begin{latin}
\begin{lstlisting}
lsb_release -a
\end{lstlisting}
\end{latin}

% --------------------------------------------------
% Results table
% --------------------------------------------------
\section{نتایج و مشاهدات}

جدول \ref{tab:commands}
کاربرد دستورات استفاده‌شده در آزمایش را نشان می‌دهد.

\begin{table}[htbp]
\centering
\caption{دستورات پایه ترمینال}
\label{tab:commands}

\begin{tabular}{cc}
\toprule
دستور & کاربرد \\
\midrule
\lr{pwd} & نمایش مسیر فعلی \\
\lr{ls -la} & نمایش جزئیات فایل‌ها \\
\lr{whoami} & نمایش نام کاربر \\
\lr{uname -a} & نمایش اطلاعات هسته \\
\lr{lsb\_release -a} & نمایش نسخه توزیع \\
\bottomrule
\end{tabular}

\end{table}

% --------------------------------------------------
% Formula example
% --------------------------------------------------
\section{نمونه فرمول ریاضی}

برای آشنایی با قابلیت فرمول‌نویسی لاتک،
میانگین سه مقدار نمونه را محاسبه می‌کنیم.

\begin{equation}
\bar{x} = \frac{x_1+x_2+x_3}{3}
\label{eq:average}
\end{equation}

اگر سه مقدار برابر با ۴، ۶ و ۸ باشند:

\[
\bar{x} = \frac{4+6+8}{3} = 6
\]

رابطه \eqref{eq:average}
فرمول میانگین حسابی را نشان می‌دهد.

% --------------------------------------------------
% Image instructions
% --------------------------------------------------
\section{اضافه کردن تصویر}

برای افزودن تصویر واقعی از مراحل آزمایش،
ابتدا پوشه‌ای با نام
\lr{images}
در کنار فایل اصلی ایجاد کنید.

سپس تصویر موردنظر را داخل آن قرار دهید.

برای مثال، پس از ذخیره تصویری با نام
\lr{terminal.png}
می‌توانید کد زیر را به سند اضافه کنید.

% Remove the comment markers after adding the image:
%
% \begin{figure}[htbp]
%     \centering
%     \includegraphics[
%         width=0.85\textwidth
%     ]{images/terminal.png}
%     \caption{خروجی ترمینال}
%     \label{fig:terminal}
% \end{figure}

% --------------------------------------------------
% Troubleshooting
% --------------------------------------------------
\section{مشکلات و راه‌حل‌ها}

در صورت بروز مشکل هنگام اجرای دستورات،
ابتدا متن خطا بررسی می‌شود.

در صورت بروز مشکل در تولید PDF نیز
تنظیمات کامپایلر، بسته‌های نصب‌شده،
فونت‌ها و مسیر فایل‌ها باید بررسی شوند.

% --------------------------------------------------
% Conclusion
% --------------------------------------------------
\section{نتیجه‌گیری}

در این آزمایش با محیط ترمینال لینوکس
و تعدادی از دستورات پایه آن آشنا شدیم.

همچنین یاد گرفتیم چگونه یک گزارش دانشگاهی
را با استفاده از لاتک تهیه کنیم و
بخش‌های مختلف آن را به‌صورت منظم نمایش دهیم.

\end{document}
